% ============================================================
\section{Limitations and Uncertainties}
\label{sec:limits}
% ============================================================

\subsection{DEM Resolution}

The 30-m Copernicus DEM GLO-30 is adequate for computing $\kT$ and
$\kM$ at most CRNS sites, where the footprint radius exceeds
\SI{100}{m}.
However, features smaller than $\sim$60 m (two pixels) are not
resolved: narrow gullies, embankments, road cuts, and individual
buildings may be missed.
In complex built or agricultural terrain, a higher-resolution LiDAR
DEM (1–5 m) is recommended.
The ray-casting algorithm accepts any DEM; swapping the 30-m tile
for a LiDAR product requires only a change in the download module.

The horizon-scan algorithm is limited to azimuths sampled at $5^\circ$
intervals (72 directions).
Narrow topographic gaps (e.g.\ a mountain pass at $4^\circ$ width)
between two blocking ridges may be missed.
Increasing $N_\phi$ reduces this error at the cost of proportionally
more CPU time.

\subsection{LULC Classification}

The ESA WorldCover 2021 product has a reported overall accuracy of
$\sim$75\% at the 10-m pixel level.
Classification errors in the heterogeneous environments common to
CRNS sites (forest edges, rice paddies, urban green space) can
propagate directly into $\kL$.
A 10\% misclassification of forest/grassland pixels corresponds to a
$\Delta\kL \approx 0.06$ error, translating to a soil-moisture bias
of approximately $\SI{0.006}{\cubic\metre\per\cubic\metre}$.

The $f_H$ values in Table~\ref{tab:fH} are literature means;
actual vegetation water content varies with phenology and rainfall.
For precise calibration, site-specific leaf water content measurements
at the time of calibration are preferable.

\subsection{Soil Properties}

SoilGrids 2.0 provides 90th-percentile uncertainty maps alongside
its predictions.
Bulk density $\rho_b$ and clay fraction have typical uncertainties
of $\pm 0.1$ g/cm$^3$ and $\pm 5$ g/kg respectively at the
250-m resolution.
These propagate into $\lw$ and $\zzz$:
\begin{align*}
  \delta\lw   &\approx 0.097\,\delta f_\mathrm{clay} \sim 0.005 \text{ g/g},\\
  \delta\zzz  &\approx \frac{8.3}{\rho_b^2 (0.0564+\thetav+\lw)}\,\delta\rho_b
               \sim \SI{2}{cm}.
\end{align*}
The SoilGrids uncertainty is a spatial mean at 250 m; real local
variability, especially in organic-rich or heterogeneous soils, can
be larger.
In-situ soil sampling at the calibration rings (Sec.~\ref{sec:methods:sampling})
remains the gold standard for $\rho_b$ and clay determination.

\subsection{Lattice Water Estimate}

Equation~\eqref{eq:lw_clay} is an empirical fit across a global
soil database.
For soils dominated by 1:1 clays (kaolinite, halloysite) the
coefficient 0.097 may overestimate $\lw$ by 20–30\%, whereas for
2:1 swelling clays (smectite, montmorillonite) the same equation
may underestimate it.
Geological characterisation (Sec.~\ref{sec:methods:sampling}) provides
a qualitative flag when the site lithology suggests an unusual
clay mineralogy.

\subsection{Atmospheric Corrections}

The ERA5-derived $\rho_{WV}$ climatology has a spatial resolution
of $\sim$30 km and may not capture local valley inversions, sea-breeze
effects, or orographic enhancement of humidity.
For sites with strong mesoscale meteorological forcing, on-site
humidity measurements are recommended for operational $f_{WV}$ correction.

The SWE estimate uses a fixed snow density $\rho_\mathrm{snow} =
\SI{250}{\kilo\gram\per\cubic\metre}$.
Fresh snow and densified spring snowpack can deviate by
$\pm\SI{100}{\kilo\gram\per\cubic\metre}$, introducing a 40\%
uncertainty in SWE.
This affects only the seasonal operability assessment (month of
$\mathrm{SWE} > \SI{150}{mm}$) and not the soil-moisture retrievals.

\subsection{Geological Data}

The Macrostrat API provides a qualitative, reconnaissance-level
geological characterisation.
At some sites — particularly in data-sparse regions (parts of Africa,
Central Asia, Antarctica) — no unit may be returned, and the
fallback is a ``not available'' flag.
Hydrogeological properties (fracturing, karst porosity, perched
water tables) cannot be retrieved and must be assessed through
field investigation.

\subsection{Geomagnetic Rigidity}

The rigidity cutoff $R_c$ from the Smart \& Shea (2019) tables
carries an uncertainty of $\sim 0.1$ GV \citep{smart2019}, which
translates to $\Delta f_{R_c} \approx 0.7\%$.
This is typically the smallest uncertainty in the $N_0$ chain and
can be neglected in most applications.

\subsection{Combined Uncertainty on \texorpdfstring{$N_0$}{N0}}

Assuming the dominant sources are independent and normally
distributed, the relative uncertainty on $N_0$ is approximately
\begin{equation}
  \frac{\sigma_{N_0}}{N_0} \approx
    \sqrt{\left(\frac{\sigma_{\kT}}{\kT}\right)^2
        + \left(\frac{\sigma_{\kM}}{\kM}\right)^2
        + \left(\frac{\sigma_{\kL}}{\kL}\right)^2
        + \left(\frac{\sigma_{\lw}}{\lw + a_2}\right)^2
    }.
  \label{eq:N0_uncertainty}
\end{equation}
For the Italian Alps example, a rough budget gives:
$\sigma_{\kT}/\kT \approx 1\%$ (DEM resolution),
$\sigma_{\kM}/\kM \approx 0.2\%$ (horizon sampling),
$\sigma_{\kL}/\kL \approx 4\%$ (LULC classification + $f_H$),
$\sigma_{\lw}/(\lw+a_2) \approx 3\%$ (SoilGrids clay uncertainty).
This gives $\sigma_{N_0}/N_0 \approx 5\%$, corresponding to a
soil-moisture uncertainty of roughly
$\SI{0.015}{\cubic\metre\per\cubic\metre}$ at
$\thetav = \SI{0.20}{\cubic\metre\per\cubic\metre}$.
This is comparable to the precision of a single calibration
measurement and underscores the value of the site characterisation
workflow.
